The Digital Farmer Assistance System addresses the critical information gap faced by small-scale farmers in rural areas by providing a centralized platform for essential agricultural data. Traditional farming operations often suffer from poor resource management due to inaccessible or delayed market trends, volatile weather patterns, and unverified crop guidance. This project presents a web-based portal developed to deliver real-time, validated crop-specific agricultural advisories, local market price updates, and a five-day localized weather forecasting mechanism integrated with extreme weather alerts. Built on a four-tier architecture featuring a Node.js runtime, an Express.js server, and a MySQL database backend, the system emphasizes scalability and low-latency access over varied bandwidths. The frontend provides a mobile-first user interface tailored to accessibility constraints. The implementation features a role-based access control paradigm separating farmer and administrative interfaces, ensuring secure and private user interactions. A core component of the system is the expert query resolution module, bridging the gap between agronomists and farmers. Functionalities including multilingual database-level support and robust offline fallback database pathways cater particularly to communities with disparate levels of digital literacy and unreliable connectivity. As a result, the platform aims to empower farmers through data-driven decisions that will simultaneously improve crop yields and promote sustained financial stability.
